% from ega4-2-fr.tex
\oldpage[IV-2]{13}
pour $g:X'\to X$ le morphisme canonique~; il est clair que $g$ est fidèlement
plat, et l'hypothèse entraîne que $g^*(\mathscr F)$ est fidèlement plat
relativement à $f\circ g$, en vertu de (2.2.11, (iii)).

\begin{corollary}[2.2.13]
\label{IV.2.2.13-fr}
\begin{enumerate}
\item[(i)] Soient $f:X\to Y$, $g:Y'\to Y$ deux morphismes,
$X'=X\times_Y Y'$, $f'=f_{(Y')}:X'\to Y'$. Si $f$ est un morphisme
fidèlement plat, il en est de même de $f'$.
\item[(ii)] Si $f:X\to X'$, $g:Y\to Y'$ sont deux $S$-morphismes fidèlement
plats,
\[
f\times_S g:X\times_S Y\longrightarrow X'\times_S Y'
\]
est fidèlement plat.
\item[(iii)] Soient $f:X\to Y$, $g:Y\to Z$ deux morphismes tels que $f$ soit
fidèlement plat. Pour que $g$ soit un morphisme plat (resp. fidèlement plat),
il faut et il suffit que $g\circ f$ soit un morphisme plat (resp. fidèlement
plat).
\end{enumerate}
\end{corollary}

\begin{proposition}[2.2.14]
\label{IV.2.2.14-fr}
Soit $f:X\to Y$ un morphisme fidèlement plat et quasi-compact. Si $X$ est
localement noethérien, il en est de même de $Y$.
\end{proposition}

La question étant locale sur $Y$, on peut supposer $Y=\operatorname{Spec}(A)$
affine~; comme $f$ est quasi-compact, il résulte de (2.2.12) que l'on peut
aussi se borner au cas où $X=\operatorname{Spec}(B)$ est affine. Alors $B$
est un anneau noethérien par hypothèse, et un $A$-module fidèlement plat
(2.2.3)~; donc $A$ est noethérien $(0_{\mathrm I},6.5.2)$.

\begin{proposition}[2.2.15]
\label{IV.2.2.15-fr}
Soit $f:X\to Y$ un morphisme fidèlement plat. Si $\mathfrak S(X)$ et
$\mathfrak S(Y)$ désignent respectivement l'ensemble des sous-préschémas de
$X$ et de $Y$, l'application $Z\mapsto f^{-1}(Z)$ de $\mathfrak S(Y)$ dans
$\mathfrak S(X)$ est injective.
\end{proposition}

Comme $f$ est surjectif, on a, pour l'ensemble sous-jacent à un
sous-préschéma $Z$ de $Y$, $f(f^{-1}(Z))=Z$. D'autre part, si $U$ est un
ouvert de $Y$ contenant $Z$ et dans lequel $Z$ est fermé, $f^{-1}(U)$ est
ouvert dans $X$, $f^{-1}(Z)$ est fermé dans $f^{-1}(U)$, et la restriction
$f^{-1}(U)\to U$ de $f$ est un morphisme fidèlement plat. On peut donc se
borner à ne considérer que des sous-préschémas \emph{fermés} de $Y$. Or, si
$Z$ est un sous-préschéma fermé de $Y$ correspondant à un Idéal
quasi-cohérent $\mathscr J$ de $\mathcal O_Y$, on sait que $f^{-1}(Z)$
correspond à l'Idéal quasi-cohérent $f^*(\mathscr J)\mathcal O_X$ de
$\mathcal O_X$ (I, 4.4.5), et comme $f$ est plat,
$f^*(\mathscr J)\mathcal O_X$ s'identifie à $f^*(\mathscr J)$. Mais
l'application $\mathscr J\mapsto f^*(\mathscr J)$ de l'ensemble des Idéaux
quasi-cohérents de $\mathcal O_Y$ dans l'ensemble des Idéaux quasi-cohérents
de $\mathcal O_X$ est injective (2.2.7), d'où la conclusion.

\begin{corollary}[2.2.16]
\label{IV.2.2.16-fr}
Soient $X$, $Y$ deux $S$-préschémas~; si $S'\to S$ est un morphisme
fidèlement plat, l'application $f\mapsto f_{(S')}$ de
$\operatorname{Hom}_S(X,Y)$ dans
$\operatorname{Hom}_{S'}(X_{(S')},Y_{(S')})$ est injective.
\end{corollary}

On a $X_{(S')}\times_{S'}Y_{(S')}=(X\times_S Y)_{(S')}$ (I, 3.3.10), donc le
morphisme projection
$X_{(S')}\times_{S'}Y_{(S')}\to X\times_S Y$ est fidèlement plat (2.2.13).
Les éléments de $\operatorname{Hom}_S(X,Y)$ correspondent biunivoquement aux
sous-préschémas de $X\times_S Y$ qui sont les \emph{graphes} de ces
$S$-morphismes (I, 5.3.11), et si $Z_f$ est le graphe de
$f\in\operatorname{Hom}_S(X,Y)$, on a $Z_{f_{(S')}}=g^{-1}(Z_f)$ (I, 5.3.12).
Il suffit donc d'appliquer à $g$ la proposition (2.2.15).

\begin{proposition}[2.2.17]
\label{IV.2.2.17-fr}
Soient $A$ un anneau, $B$ une $A$-algèbre telle que $B$ soit un $A$-module
fidèlement plat et de présentation finie. Alors l'homomorphisme structural
$\varphi:A\to B$ est un isomorphisme du $A$-module $A$ sur un \emph{facteur
direct} du $A$-module $B$. Si $A$ est un anneau local, $B$ est un $A$-module
libre et il existe une base de ce module contenant l'élément unité de $B$.
\end{proposition}

